\chapter{Complex Numbers and Euler's Identity}\label{ch:cx}

$e^{i\pi} + 1 = 0$ is the identity on the cover of this book's notebook, and for most of the
project the engine could not compute it, because it had no $i$. Adding one is not a matter of a
symbol. Every theorem in the preceding chapters is about real numbers; a term containing $i$ has no
real meaning, and the rules that are safe over $\mathbb{R}$ are not all safe over $\mathbb{C}$.
This chapter relearns the complex numbers from their construction, proves Euler's formula from the
exponential series, and then meets the engine's answer to the problem of meaning: a third
semantics, and a ledger with two columns.

\section{The construction}

\begin{definition}
  $\mathbb{C}$ is the set $\mathbb{R}^2$ with $(a, b) + (c, d) = (a + c, b + d)$ and
  $(a, b)(c, d) = (ac - bd, ad + bc)$. Write $a + bi$ for $(a, b)$, so that $i = (0, 1)$ and
  $i^2 = (-1, 0) = -1$.
\end{definition}

\begin{theorem}
  $\mathbb{C}$ is a field.
\end{theorem}
\begin{proof}
  The ring axioms are coordinate computations, each a polynomial identity in four or six real
  variables; associativity of multiplication is the longest and is still one line. The inverse of
  $z = a + bi \neq 0$ is
  \[
    z^{-1} = \frac{a - bi}{a^2 + b^2},
  \]
  since $(a + bi)(a - bi) = a^2 + b^2 > 0$; that is the only place the field structure of
  $\mathbb{R}$ beyond a ring is used.
\end{proof}

The equivalent description, which explains why the rules look the way they do, is that $\mathbb{C}
= \mathbb{R}[x]/(x^2 + 1)$: polynomials in a formal $x$, reduced by the single relation $x^2 = -1$.
Every polynomial in $i$ collapses to a linear one, $a + bi$, and that collapse is what the engine's
arithmetic rule performs.

\begin{definition}
  The conjugate of $z = a + bi$ is $\bar z = a - bi$; the modulus is $|z| = \sqrt{a^2 + b^2} =
  \sqrt{z\bar z}$; the real and imaginary parts are $\operatorname{Re} z = a = \frac{z + \bar
  z}{2}$ and $\operatorname{Im} z = b = \frac{z - \bar z}{2i}$.
\end{definition}

Conjugation is a ring automorphism, $\overline{zw} = \bar z \bar w$ and $\overline{z + w} = \bar z
+ \bar w$, because it is the other root of $x^2 + 1$; and $|zw| = |z||w|$ follows by taking
$\sqrt{\cdot}$ of $zw\,\overline{zw} = z\bar z\, w \bar w$.

\section{The exponential, once more}

The series of Chapter~\ref{ch:explog} converges for every complex $z$ by the same ratio test, and
the functional equation $\exp(z + w) = \exp z \exp w$ has the same proof. Splitting the series of
$\exp(i\theta)$ for real $\theta$ into even and odd powers of $i$ gives

\begin{theorem}[Euler's formula]
  For real $\theta$, $\exp(i\theta) = \cos\theta + i \sin\theta$.
\end{theorem}
\begin{proof}
  $i^n$ cycles through $1, i, -1, -i$, so
  \[
    \sum_{n \ge 0} \frac{(i\theta)^n}{n!}
      = \sum_{k \ge 0} (-1)^k \frac{\theta^{2k}}{(2k)!} + i \sum_{k \ge 0} (-1)^k \frac{\theta^{2k+1}}{(2k+1)!}
      = \cos\theta + i\sin\theta,
  \]
  the rearrangement being justified by absolute convergence, and the two series being the
  definitions of $\cos$ and $\sin$.
\end{proof}

With $\theta = \pi$ and the values $\cos\pi = -1$, $\sin\pi = 0$: $e^{i\pi} = -1$. The identity is
two facts, then --- Euler's formula, and the exact values of the trigonometric functions at $\pi$
--- and the engine needs both as rules.

\section{Exact values on the unit circle}

$\sin$ and $\cos$ at $0, \pi/6, \pi/4, \pi/3, \pi/2$ are the five values every student memorizes:
$\sin$ runs $0, \frac{1}{2}, \frac{\sqrt 2}{2}, \frac{\sqrt 3}{2}, 1$ and $\cos$ runs them
backwards. Everything else at a rational multiple of $\pi$ follows from four symmetries:
periodicity $\sin(\theta + 2\pi) = \sin\theta$, the half-turn $\sin(\theta + \pi) = -\sin\theta$,
$\cos(\theta + \pi) = -\cos\theta$, and the reflection $\sin(\pi - \theta) = \sin\theta$,
$\cos(\pi - \theta) = -\cos\theta$. Given $q\pi$, reduce $q$ modulo $2$, then by the half-turn into
$[0, 1)$, then by the reflection into $[0, \frac{1}{2}]$, and look up. The engine's rule
\rn{cx.exact-trig} does exactly this reduction on the rational $q$, and its proof over
$\mathbb{R}$ is the four Mathlib identities (\lc{Real.sin_add_int_mul_two_pi},
\lc{Real.sin_add_pi}, \lc{Real.sin_pi_sub} and their cosine forms) followed by the five table
entries (\lc{Real.sin_pi_div_six} and friends). $\tan$ is the quotient where the cosine is not
zero.

\section{Why the ordering does not object}

Euler's formula in general, $\exp(i\theta) \to \cos\theta + i\sin\theta$, duplicates $\theta$, and
Chapter~\ref{ch:wf} explained why no weight of the engine's shape can pay for a duplication. The
rule \rn{cx.euler} therefore fires only where both values are \emph{exact}: the result then
contains no $\theta$ at all, and the exponential node it replaces outweighs any Gaussian numeral.
This is also what Mathematica does: \texttt{Exp[I x]} stays as it is, \texttt{Exp[I Pi]} becomes
$-1$.

The same accounting shapes the arithmetic. $i^2 \to -1$, a product of two Gaussian numerals, an
integer power, a conjugate, a modulus --- each replaces a node by a smaller Gaussian numeral. One
case is not affordable: $(1 + i)^{-1} \to \frac{1}{2} - \frac{1}{2} i$ weighs more after than
before, because non-integer numerals are heavy (Chapter~\ref{ch:cannot} explains why they must
be). But $1/(1 + i)$ arrives as the product $1 \cdot (1 + i)^{-1}$, and the product with its
leading $1$ still on it is heavy enough; so the arithmetic rule treats an integer power of a
Gaussian numeral as a Gaussian factor of a product, and runs before the identity rule that would
strip the $1$. Rule order inside the pipeline is not part of the termination proof, which is per
rule, so this costs nothing.

\section{A third semantics}

\begin{minted}{lean}
def evalC (ρ : EnvC) : Expr → ℂ
  | .num q => (q.val : ℂ)
  | .var x => ρ x
  | .add es => sumC ρ es
  | .mul es => prodC ρ es
  | .pow b e => (evalC ρ b) ^ (evalC ρ e)
  | .fn f [] => constC f          -- π ↦ Real.pi, i ↦ Complex.I
  | .fn f [a] => applyFnC f (evalC ρ a)
  | .fn _ _ => 0
  | .matrix _ => 0
\end{minted}

Powers are Mathlib's \lc{Complex.cpow}, the principal branch: $z^w = \exp(w \log z)$ with the
logarithm's imaginary part in $(-\pi, \pi]$. That choice has consequences a real-number rule never
met. $(e^a)^b = e^{ab}$ is false in general --- $(e^{2\pi i})^{1/2} = 1^{1/2} = 1$ while $e^{\pi i}
= -1$ --- so the rule $\ln(e^x) \to x$ of Chapter~\ref{ch:explog}, unconditional over $\mathbb{R}$
for the exponential's argument, is refuted over $\mathbb{C}$ at $x = 2\pi i$. And $b^m b^n =
b^{m+n}$, which needed $b > 0$ over $\mathbb{R}$, needs only $b \neq 0$ over $\mathbb{C}$
(\lc{Complex.cpow_add}): a weaker condition in a bigger field, because the principal branch is
chosen consistently.

So the ledger acquires a second column. Every rule has its status over $\mathbb{R}$ as before, and
a rule with a theorem over $\mathbb{C}$ has a second status there. The four algebraic rules of the
simplifier and the structural power rules are verified in both. Collecting powers is conditional
in both, with different conditions. The function rule is conditional in both, for different
reasons. The complex rules of this chapter are verified over $\mathbb{C}$ and have no status over
$\mathbb{R}$, where $i$ means nothing. And a cell whose input or output mentions $i$ is read in
the complex semantics: the notebook shows, for every step, the rule's status \emph{there}, and a
rule proved only over $\mathbb{R}$ shows as unverified in such a cell, with a note saying so.
That is the honest state, and it is the cost of adding a number.

\begin{minted}{lean}
theorem euler_soundC (ρ : EnvC) {e : Expr} {res : RuleResult}
    (h : euler.apply e = some res) : evalC ρ res.result = evalC ρ e
theorem exactTrig_soundR (ρ : EnvR) {e : Expr} {res : RuleResult}
    (h : exactTrig.apply e = some res) : evalR ρ res.result = evalR ρ e
theorem not_functionRules_soundC : ¬ RuleSoundC functionRules
\end{minted}

\begin{trybox}
\begin{minted}{text}
ℯ^(pi*i)
exp(i*pi/4)
(1+i)*(2-i)
1/(1+i)
abs(3+4i)
cos(7pi/6)
\end{minted}
  Type \texttt{\textbackslash e} and \texttt{\textbackslash pi} for the constants. Open the work
  for the first: the power becomes an exponential, then Euler's formula fires because both values
  are exact. Click a step and see "Proof status over ℂ".
\end{trybox}

\section{Exercises}

\begin{exercisebox}[The quotient]
  $\mathbb{C} = \mathbb{R}[x]/(x^2+1)$: show that every class has a unique representative of
  degree at most $1$, and that the multiplication of representatives is the one in the definition.
\end{exercisebox}

\begin{exercisebox}[Principal values]
  Compute $i^i$ with the principal branch, and find two complex $a, b$ for which
  $(e^a)^b \neq e^{ab}$ with $|a|$ as small as you can make it.
\end{exercisebox}

\begin{exercisebox}[The table]
  Prove $\sin\frac{\pi}{6} = \frac{1}{2}$ from $\sin 3\theta = 3\sin\theta - 4\sin^3\theta$, and
  $\cos\frac{\pi}{4} = \frac{\sqrt 2}{2}$ from $\cos 2\theta = 2\cos^2\theta - 1$.
\end{exercisebox}

\begin{exercisebox}[Where the ordering stops]
  Show that with the engine's weights $(1 + i)^{-1}$ weighs less than $\frac{1}{2} - \frac{1}{2}i$,
  and that $1 \cdot (1 + i)^{-1}$ weighs more. Which weight would have to change to make the bare
  inverse a rule, and what would that change break?
\end{exercisebox}
